% !TEX program = xelatex
% ============================================================
%  مرجع واجهة البرمجة — API Reference Template
%  قالب لوثائق مرجع واجهة برمجة التطبيقات (API Reference).
%  Compile with: xelatex main.tex (run twice)
% ============================================================
%  Features:
%    - Authentication overview, base URL
%    - Endpoints table of contents
%    - Per-endpoint sections: parameters, request/response examples
%    - Code blocks wrapped in english environment (LTR-safe)
%    - Error codes table per endpoint
% ============================================================

\documentclass[11pt,a4paper]{article}

\usepackage[a4paper,margin=2.5cm,top=2.5cm,bottom=2.5cm,headheight=16pt]{geometry}
\usepackage{graphicx}
\usepackage{array}
\usepackage{booktabs}
\usepackage{tabularx}
\usepackage{multirow}
\usepackage{enumitem}
\usepackage{float}
\usepackage{caption}
\usepackage{titlesec}
\usepackage{fancyhdr}
\usepackage{setspace}
\usepackage{xcolor}
\usepackage{amsmath}
\usepackage{amssymb}
\usepackage{listings}

\usepackage{bidi}
\usepackage[no-math]{fontspec}
\usepackage[quiet]{polyglossia}

\setmainfont[Scale=1]{Amiri-Regular.ttf}
\newfontfamily\arabicfont[Script=Arabic,Scale=0.95]{Amiri-Regular.ttf}
\newfontfamily\englishfont[Scale=0.95]{TeX Gyre Termes}
\setmonofont{TeX Gyre Cursor}

\setdefaultlanguage[calendar=gregorian,hijricorrection=1]{arabic}
\setotherlanguage[variant=british]{english}

\usepackage[breaklinks,hidelinks]{hyperref}
\hypersetup{pdftitle={مرجع واجهة البرمجة},pdfauthor={اسم الشركة}}

\addto\captionsarabic{%
  \renewcommand{\contentsname}{الفهرس}%
  \renewcommand{\figurename}{الشكل}%
  \renewcommand{\tablename}{الجدول}%
}

\titleformat{\section}{\Large\bfseries}{\thesection}{0.7em}{}
\titleformat{\subsection}{\large\bfseries}{\thesubsection}{0.6em}{}
\titlespacing*{\section}{0pt}{1.4ex plus 0.4ex}{0.9ex plus 0.3ex}

\pagestyle{fancy}
\fancyhf{}
\fancyhead[R]{\small\itshape مرجع واجهة البرمجة (API Reference)}
\fancyhead[L]{\small اسم المنتج}
\fancyfoot[C]{\small\thepage}
\renewcommand{\headrulewidth}{0.3pt}

\setlist[itemize]{topsep=0.3em, itemsep=0.15em, parsep=0pt}
\setlist[enumerate]{topsep=0.3em, itemsep=0.15em, parsep=0pt}

\newcolumntype{L}[1]{>{\raggedright\arraybackslash}p{#1}}
\newcolumntype{C}[1]{>{\centering\arraybackslash}p{#1}}
\newcolumntype{R}[1]{>{\raggedleft\arraybackslash}p{#1}}

\definecolor{headerColor}{HTML}{1A3C6B}
\definecolor{codeBg}{HTML}{F7F7F7}
\definecolor{codeFrame}{HTML}{CCCCCC}
\definecolor{methodGet}{HTML}{2E7D32}
\definecolor{methodPost}{HTML}{1565C0}

\lstset{
  basicstyle=\ttfamily\small,
  breaklines=true,
  frame=single,
  rulecolor=\color{codeFrame},
  backgroundcolor=\color{codeBg},
  showstringspaces=false,
  columns=fullflexible,
  keepspaces=true,
}

\setstretch{1.3}

\begin{document}

% ============================================================
% TITLE BLOCK
% ============================================================
\begin{center}
{\LARGE\bfseries مرجع واجهة البرمجة (API Reference)}\\[0.5em]
{\Large اسم المنتج --- وثائق الـ \LR{API}}\\[1em]
\end{center}

% ---- Version and status table ----
\begin{table}[H]
\centering
\renewcommand{\arraystretch}{1.4}
\begin{tabularx}{\textwidth}{L{3.5cm} X L{3.5cm} L{3cm}}
\toprule
\textbf{اسم المنتج:} & \multicolumn{3}{l}{---} \\
\textbf{الإصدار:} & \LR{v1.0.0} & \textbf{الحالة:} & مستقر (Stable) \\
\textbf{اسم الشركة:} & --- & \textbf{التاريخ:} & --- \\
\bottomrule
\end{tabularx}
\end{table}

\vspace{0.5em}

% ============================================================
% 1. AUTHENTICATION
% ============================================================
\section{المصادقة (Authentication)}
\label{sec:auth}

% TODO: صف آلية المصادقة المستخدمة.
تستخدم واجهة البرمجة مصادقة عبر رمز حامل (\LR{Bearer Token}). يجب إرفاق الرمز
في ترويسة الطلب (\LR{Authorization header}) كما يلي:

\begin{english}
\begin{lstlisting}[language=bash]
Authorization: Bearer <YOUR_API_TOKEN>
\end{lstlisting}
\end{english}

\begin{itemize}
    \item تُصدر الرموز من لوحة التحكم (\LR{Dashboard}) الخاصة بالمطور.
    \item تنتهي صلاحية الرمز بعد \LR{24} ساعة من إصداره.
    \item في حال انتهاء الصلاحية يُعاد الطلب مع رمز الحالة \LR{401 Unauthorized}.
\end{itemize}

% ============================================================
% 2. BASE URL
% ============================================================
\section{عنوان الخادم الأساسي (Base URL)}
\label{sec:base-url}

% TODO: حدد عنوان الخادم الأساسي لكل بيئة.
\begin{table}[H]
\centering
\renewcommand{\arraystretch}{1.3}
\begin{tabularx}{\textwidth}{L{3cm} X}
\toprule
\textbf{البيئة (Environment)} & \textbf{العنوان (URL)} \\
\midrule
الإنتاج (\LR{Production}) & \LR{https://api.example.com/v1} \\
الاختبار (\LR{Staging})   & \LR{https://staging.api.example.com/v1} \\
التطوير (\LR{Development}) & \LR{https://dev.api.example.com/v1} \\
\bottomrule
\end{tabularx}
\end{table}

% ============================================================
% 3. ENDPOINTS TOC
% ============================================================
\section{فهرس نقاط النهاية (Endpoints)}
\label{sec:endpoints-toc}

% TODO: حدّث قائمة نقاط النهاية.
\begin{table}[H]
\centering
\renewcommand{\arraystretch}{1.3}
\begin{tabularx}{\textwidth}{C{1.5cm} L{4.5cm} C{2cm} X}
\toprule
\textbf{القسم} & \textbf{المسار} & \textbf{الطريقة} & \textbf{الوصف} \\
\midrule
\ref{sec:ep-get-users} & \LR{/users}     & \LR{GET}  & استرجاع قائمة المستخدمين. \\
\ref{sec:ep-post-users} & \LR{/users}    & \LR{POST} & إنشاء مستخدم جديد. \\
\bottomrule
\end{tabularx}
\end{table}

% ============================================================
% 4. ENDPOINT: GET /users
% ============================================================
\section{استرجاع المستخدمين: \LR{GET /users}}
\label{sec:ep-get-users}

\subsection{المعاملات (Parameters)}
% TODO: املأ معاملات الطلب.
\begin{table}[H]
\centering
\renewcommand{\arraystretch}{1.3}
\begin{tabularx}{\textwidth}{L{3cm} C{2cm} C{2cm} X}
\toprule
\textbf{المعامل} & \textbf{النوع} & \textbf{إلزامي} & \textbf{الوصف} \\
\midrule
\LR{page}    & \LR{integer} & لا  & رقم الصفحة (افتراضي \LR{1}). \\
\LR{limit}   & \LR{integer} & لا  & عدد النتائج لكل صفحة (افتراضي \LR{20}). \\
\LR{search}  & \LR{string}  & لا  & نص البحث في الاسم والبريد. \\
\bottomrule
\end{tabularx}
\end{table}

\subsection{مثال الطلب (Request Example)}
\begin{english}
\begin{lstlisting}[language=bash]
curl -X GET "https://api.example.com/v1/users?page=1&limit=20" \
  -H "Authorization: Bearer <YOUR_API_TOKEN>"
\end{lstlisting}
\end{english}

\subsection{مثال الاستجابة (Response Example)}
\begin{english}
\begin{lstlisting}[]
{
  "data": [
    {
      "id": 1,
      "name": "Ahmed Ali",
      "email": "ahmed@example.com",
      "role": "admin"
    }
  ],
  "page": 1,
  "total": 1
}
\end{lstlisting}
\end{english}

\subsection{رموز الأخطاء (Error Codes)}
\begin{table}[H]
\centering
\renewcommand{\arraystretch}{1.3}
\begin{tabularx}{\textwidth}{C{2cm} L{4cm} X}
\toprule
\textbf{الرمز} & \textbf{الاسم} & \textbf{الوصف} \\
\midrule
\LR{200} & \LR{OK}                & تم الاسترجاع بنجاح. \\
\LR{401} & \LR{Unauthorized}      & رمز المصادقة غير صالح أو مفقود. \\
\LR{403} & \LR{Forbidden}         & لا تملك صلاحية الوصول. \\
\LR{429} & \LR{Too Many Requests} & تجاوز حد معدل الطلبات. \\
\bottomrule
\end{tabularx}
\end{table}

% ============================================================
% 5. ENDPOINT: POST /users
% ============================================================
\section{إنشاء مستخدم: \LR{POST /users}}
\label{sec:ep-post-users}

\subsection{المعاملات (Parameters)}
% TODO: املأ معاملات الطلب.
\begin{table}[H]
\centering
\renewcommand{\arraystretch}{1.3}
\begin{tabularx}{\textwidth}{L{3cm} C{2cm} C{2cm} X}
\toprule
\textbf{المعامل} & \textbf{النوع} & \textbf{إلزامي} & \textbf{الوصف} \\
\midrule
\LR{name}  & \LR{string}  & نعم & اسم المستخدم الكامل. \\
\LR{email} & \LR{string}  & نعم & البريد الإلكتروني (فريد). \\
\LR{role}  & \LR{string}  & لا  & الدور (افتراضي \LR{user}). \\
\bottomrule
\end{tabularx}
\end{table}

\subsection{مثال الطلب (Request Example)}
\begin{english}
\begin{lstlisting}[]
curl -X POST "https://api.example.com/v1/users" \
  -H "Authorization: Bearer <YOUR_API_TOKEN>" \
  -H "Content-Type: application/json" \
  -d '{
    "name": "Ahmed Ali",
    "email": "ahmed@example.com",
    "role": "admin"
  }'
\end{lstlisting}
\end{english}

\subsection{مثال الاستجابة (Response Example)}
\begin{english}
\begin{lstlisting}[]
{
  "data": {
    "id": 1,
    "name": "Ahmed Ali",
    "email": "ahmed@example.com",
    "role": "admin",
    "created_at": "2025-01-15T10:30:00Z"
  }
}
\end{lstlisting}
\end{english}

\subsection{رموز الأخطاء (Error Codes)}
\begin{table}[H]
\centering
\renewcommand{\arraystretch}{1.3}
\begin{tabularx}{\textwidth}{C{2cm} L{4cm} X}
\toprule
\textbf{الرمز} & \textbf{الاسم} & \textbf{الوصف} \\
\midrule
\LR{201} & \LR{Created}            & تم إنشاء المستخدم بنجاح. \\
\LR{400} & \LR{Bad Request}        & بيانات الطلب غير صالحة. \\
\LR{409} & \LR{Conflict}           & البريد الإلكتروني مستخدم بالفعل. \\
\LR{401} & \LR{Unauthorized}       & رمز المصادقة غير صالح. \\
\bottomrule
\end{tabularx}
\end{table}

\end{document}
